% Source: https://artofproblemsolving.com/community/c3h441901_help
% Stable identity: aops-topic-441901-q04
Angle $RQP$ is the right angle between the positive axes, so its bisector is the line $y=x$:

\begin{asy}
usepackage("amsmath");
size(6cm);
defaultpen(fontsize(10pt));
pair P=(0,5),Q=(0,0),R=(10,0),Sp=(10/3,10/3);
draw(P--Q--R--cycle,linewidth(1));
draw(Q--Sp,blue+linewidth(0.9));
real s=0.45; draw((s,0)--(s,s)--(0,s));
draw(arc(Q,1.3,0,45),blue);draw(arc(Q,1.5,45,90),blue);
dot(Sp);
label("$P(0,5)$",P,W);label("$Q(0,0)$",Q,SW);label("$R(10,0)$",R,SE);label("$S$",Sp,NE);
label("$45^\circ$",Q+2.3*dir(22.5),blue);label("$45^\circ$",Q+2.5*dir(69),blue);
label("\mbox{}",(min(currentpicture,user=true)+max(currentpicture,user=true))/2,black); // keeps the text after the figure black when LaTeX typesets the labels inline
\end{asy}

The line $PR$ has intercepts $10$ and $5$, so its equation is $\dfrac {x}{10}+\dfrac {y}{5}=1$. Putting $y=x$ gives $x\left(\dfrac{1}{10}+\dfrac{1}{5}\right)=1$, so $x=\dfrac{10}{3}$ and $S=\left(\dfrac{10}{3},\dfrac{10}{3}\right)$. Hence
\[
QS=\dfrac{10}{3}\sqrt2=\boxed{\dfrac{10\sqrt2}{3}}.
\]
(By the angle bisector theorem, $\dfrac{PS}{SR}=\dfrac{QP}{QR}=\dfrac{1}{2}$, which gives the same point $S$.)
